% Раздел о происхождении результатов, воспроизводимости и роли ИИ.
\section{Происхождение результатов, воспроизводимость и роль ИИ}
\label{sec:origin}

Работа сочетает доказательства, точные компьютерные сертификаты и обширный
численный поиск. Чтобы читатель не восстанавливал статус каждого утверждения по
контексту, ниже явно описаны шкала статусов, состав программного обеспечения,
проведённые независимые проверки и использование систем искусственного
интеллекта.

\subsection{Статусы утверждений}

Шкала \textbf{[Т]}\,/\,\textbf{[С]}\,/\,\textbf{[Ч]}\,/\,\textbf{[Э]} (с
модификатором \textbf{[Ч]}${}^{*}$ для кандидатов без сертификата диаметра)
введена в \S\ref{sec:intro}; уточним, что к чему относится. Статус \textbf{[Т]} имеют
предложение~\ref{prop:crit}, леммы~\ref{lem:sym} и~\ref{lem:mid}, планарная
теорема (предложение~\ref{prop:planar}), правило вещественного множителя
(предложение~\ref{prop:mult}), инрадиусная лемма~\ref{lem:inradius} со
следствием~\ref{cor:no-subgroup}, обе леммы о ламинировании
(\ref{lem:P1} и~\ref{lem:P2}), оба
экрана индекса (предложения~\ref{prop:mink} и~\ref{prop:inrfloor}),
продуктовое исчисление (предложение~\ref{prop:product}) вместе с
теоремой~\ref{thm:dim10} об оценке $45619$ и её следствиями для $\R^{25}$
и~$\R^{26}$, и пол мозаичного
принципа (теорема~\ref{thm:floor}). Равенство в тождестве~\eqref{eq:eis} для
конкретных решёток ($A_2$, $D_4$, $E_6^*$, $E_8$) установлено \emph{точными
вычислениями} и потому имеет статус \textbf{[С]}; нижняя граница <<$\ge$>>
доказана для всех эйзенштейновых решёток и имеет статус \textbf{[Т]}. Статус
\textbf{[С]} --- теоремы~\ref{thm:main},~\ref{thm:r5},~\ref{thm:r7} вместе с
сертификатом \S\ref{sec:comp}, сопутствующие конструкции индексов $46$ и~$48$
(\S\ref{sec:wide}), мозаичные ступени \S\ref{ssec:quant} и
интервальная ступень $\R^3/15$ (\S\ref{sec:stair}): каждая
сведена к конечному списку неравенств между явно выписанными дробями, так что
результат не зависит от корректности оптимизирующего кода --- достаточно
доверять программе, проверяющей этот список. Статус \textbf{[Ч]} имеют все
<<численные оптимумы>> в таблицах и оценки, опирающиеся на кусочные
сертификаты диаметра ($1029$, $7203$, $9604$, $28812$ и ширины ламинирований:
плавающая точка сведена к значениям явных квадратик в вершинах явных
политопов, но рациональная запись ещё не выполнена); модификатор
\textbf{[Ч]}${}^{*}$ несёт кандидат $21609$ (\S\ref{ssec:dim10lam}), чей
сертификат формально проходит, но с запасом на грани разрешения плавающей
точки. Статус \textbf{[Э]} ---
все отрицательные экраны: \S\ref{ssec:dim912neg}, \S\ref{sec:open} и
приложение~\ref{app:screens}.

Насколько существенна разница между \textbf{[С]} и \textbf{[Э]}, показывает
история индекса~$134$ в \S\ref{app:d5}: тамошний барьер был обойдён не
уточнением той же формы, а переходом к другой. Оценки $45$, $132$ и $1323$ не
опираются на \textbf{[Ч]} и \textbf{[Э]}; оценки $1029$, $7203$ и $28812$
имеют статус \textbf{[Ч]} указанной повышенной пробы.

\subsection{Программное обеспечение и данные}

Весь код и все промежуточные данные открыты:
\begin{center}
\url{https://github.com/ivaleo/Chromatic}\quad (лицензия MIT)
\end{center}

Репозиторий состоит из трёх пакетов и каталога данных:
\begin{itemize}[leftmargin=1.4em,itemsep=1pt]
\item \texttt{voronoi4d} (каталог \texttt{voronoi/}) --- эталонная реализация на
Python: построение ячейки Вороного, расстояние точка\,--\,многогранник каскадом
проекций, точная рациональная сертификация;
\item \texttt{combigeo} --- быстрое ядро на C\raise.1ex\hbox{\small++}
(перечисление коротких векторов методом Финке--Похста, GJK, локальный поиск
\texttt{min\_conflicts}, безвершинный радиус покрытия) с привязками к Python;
\item \texttt{chromatic} --- над-пакет, выполняющий перекрёстную сверку двух
независимых бэкендов (\texttt{compare\_backends});
\item \texttt{audit-data/} --- сертификаты, протоколы независимых аудитов и
результаты всех кампаний (${\sim}220$ файлов JSON, около $60$\,МБ), а также
\texttt{audit-data/chromatic\_research/} --- $125$ скриптов поисковых кампаний
в размерностях $2$--$26$, которыми получены результаты \S\ref{sec:big},
\S\ref{sec:dim912}--\ref{sec:tilings} и \S\ref{sec:open}.
\end{itemize}

Совокупный объём --- около $51\,000$ строк Python и $3\,400$ строк
C\raise.1ex\hbox{\small++}. Каждый сертификат содержит данные, достаточные для
независимой перепроверки без обращения к коду проекта: рациональную форму Грама,
матрицу подрешётки, точный квадрат радиуса покрытия, полный список
потенциально опасных векторов и KKT-сертификат проекции для каждого из них.

\subsection{Независимые проверки}

Ключевые величины проверялись не одним, а несколькими несовпадающими путями.

\begin{itemize}[leftmargin=1.4em,itemsep=2pt]
\item \emph{Две независимые реализации.} Расстояние
$\dist(v/2,V_0)$ вычисляется каскадом ортогональных проекций
(\texttt{voronoi4d}), алгоритмом GJK и решением задачи квадратичного
программирования (\texttt{combigeo}); на всех тестах три способа совпадают до
девяти значащих цифр.

\item \emph{Аудит без Qhull.} Для конструкции индекса~$132$ ячейка Вороного
строится повторно программой, не использующей ни Qhull, ни вещественный
перечислитель коротких векторов: она заново находит $31$ ненулевой
parity-класс, $62$ фасеты, обходит граф из $720$ вершин и $1800$ рёбер и
воспроизводит те же $R^2$, $D_{\min}^2$ и все KKT-сертификаты.

\item \emph{Устойчивость к рационализации.} Та же численная форма
рационализована с разными знаменателями ($10^4$, $10^5$, $10^6$): отношение
меняется в четвёртом знаке после запятой ($1{,}010771$, $1{,}010898$,
$1{,}010905$), и во всех трёх случаях строгий запас для $\ell=1{,}01$
сохраняется.

\item \emph{Повторная сертификация с нуля.} Все три \textbf{[С]}-сертификата
(индексы $45$, $132$ и $1323$) были заново воспроизведены программами, не
разделяющими с основным конвейером ни строчки кода: независимо найдены релевантные векторы, перечислены вершины ячейки
($120$, $720$ и $39\,296$), вычислен точный радиус покрытия и проверены все
KKT-сертификаты проекций; итоговые дроби совпали посимвольно. Эта
перепроверка выполнена в отдельной сессии с ИИ-ассистентом (аудит от
05.08.2026, протокол --- \texttt{journal/AUDIT-2026-08-05.md}), поэтому
независимость здесь означает независимость \emph{кода}, а не независимость
исполнителя.

\item \emph{Контроль геометрии.} Корректность построенной ячейки проверяется
соотношением Эйлера для $f$-вектора, центральной симметрией множества вершин и
равенством объёма определителю решётки. Полнота перечисления вершин
дополнительно контролируется тождеством
$\sum_{\text{вершины}}|\det(v_1,\dots,v_n)|=(n+1)!$, где $v_i$ --- релевантные
векторы активных фасет: сумма нормированных объёмов симплексов Делоне,
инцидентных узлу решётки, не зависит от комбинаторного типа ячейки.
\end{itemize}

\subsection{Использование систем искусственного интеллекта}

Работа выполнена с систематическим использованием ИИ-ассистента: вычислительный
конвейер, скрипты поисковых кампаний, иллюстрации и первоначальные редакции
текста статьи создавались в диалоговых сессиях с ним. Использовались
современные большие языковые модели общего назначения в агентной оболочке для
работы с кодом; конкретные версии менялись за время работы. Это зафиксировано
в истории репозитория: все коммиты, кроме двух ранних (\texttt{ad9279c} и
\texttt{9307135}), созданных тем же способом, несут трейлеры
\texttt{Co-Authored-By} с указанием использованной модели и идентификатора
сессии.

Конкретно с помощью ИИ-ассистента выполнены:
\begin{itemize}[leftmargin=1.4em,itemsep=1pt]
\item реализация всех алгоритмов (\S\ref{sec:algo}), включая ядро на
C\raise.1ex\hbox{\small++} и точные рациональные сертификаторы;
\item планирование, запуск и сопровождение длительных вычислительных кампаний,
в том числе исчерпывающих переборов объёмом до $2\cdot10^9$ конфигураций;
\item построение иллюстраций из данных экспериментов;
\item подготовка редакций текста статьи и сопроводительной документации.
\end{itemize}

Математическая рамка работы принадлежит не ИИ: метод раскраски по решёткам
Вороного и лемма о середине восходят к Иванову~\cite{Ivanov2011}, конструкции и
оценки-ориентиры $49$, $140$, $343$, $1372$, $2401$, $17\,253$ --- к Арману,
Бондаренко, Прымаку и Радченко~\cite{ABPR}, глобальный оптимизатор Q-поиска ---
к~А.~Горнову~\cite{GornovQ} (порт выполнен с исходного текста
первоисточника). Разделение вкладов внутри самого диалога таково: постановки
задач, направление поиска, критерии отбора и итоговая проверка результатов
принадлежат автору; конкретные механизмы --- выход в решётки общего
положения, CSP-переформулировка, примарный Radon-поиск, чередование ядра и
формы Грама, ламинирование, продуктовое исчисление ширин --- возникали в ходе
диалога с ассистентом и принимались или отбрасывались автором.

\paragraph{Что это меняет для читателя.} Для математической проверяемости
теорем --- ничего, и в этом смысл выбранной схемы работы: происхождение
поискового кода несущественно после того, как полный сертификат независимо
проверен. Оценки со статусом \textbf{[С]} сведены к конечным спискам
неравенств между явно выписанными рациональными числами; такой список
проверяется программой в несколько десятков строк, не имеющей ничего общего с
кодом, которым он получен, и его корректность не зависит ни от надёжности
оптимизирующего кода, ни от того, кем или чем этот код написан. Численные
результаты \textbf{[Ч]} и отрицательные экраны \textbf{[Э]} такой гарантии не
имеют и приводятся именно как эвристические ориентиры.

Ответственность за все утверждения, включая формулировки, полученные с помощью
ИИ-ассистента, несёт автор.
